\documentclass[11pt]{article}

\usepackage[margin=1in]{geometry}
\usepackage{amsmath,amssymb,amsfonts,amsthm}
\usepackage{booktabs}
\usepackage{hyperref}
\usepackage{listings}
\usepackage{xcolor}
\usepackage{enumitem}

\hypersetup{colorlinks=true,linkcolor=blue,urlcolor=blue,citecolor=blue}

\lstset{
  basicstyle=\ttfamily\small,
  breaklines=true,
  columns=fullflexible,
  frame=single,
  showstringspaces=false
}

\newtheorem{theorem}{Theorem}
\newtheorem{proposition}{Proposition}
\newtheorem{remark}{Remark}

\title{Squeezed-Fock State Preparation for CV-DV LCHS:\\
Parameter Analysis and Practical Limitations\\[4pt]
\large Technical Note for \texttt{clacod\_heat1d\_qutip.py}}
\author{}
\date{\today}

\begin{document}
\maketitle

\begin{abstract}
We document and analyze the squeezed-Fock state preparation method for the
hybrid continuous-variable/discrete-variable (CV-DV) linear combination of
Hamiltonian simulation (LCHS), applied to the 1D heat equation with four
interior grid points (two DV qubits).
Through systematic numerical experiments we identify the mathematical
constraints on the squeezing parameters $(r, r')$, the kernel parameter
$\beta$, and the position-operator convention, and explain why these choices
are tightly coupled.
The key finding is a tension between two competing requirements: large $r$
is needed for $\psi_r(x) \in L^2(\mathbb{R})$, but large $r$ also pushes the
squeezed vacuum beyond any practical Fock truncation.
We report the best achievable accuracy and compare the squeezed-Fock method
against the position-eigenbasis alternative.
\end{abstract}

%% ====================================================================
\section{Problem Statement}
%% ====================================================================

We consider the 1D heat equation with homogeneous Dirichlet boundary conditions,
discretized to four interior points:
\begin{equation}\label{eq:ode}
  \frac{d}{dt}\mathbf{u}(t) = -A\,\mathbf{u}(t), \qquad
  A = \frac{\alpha}{h^2}
  \begin{bmatrix}
    2 & -1 & 0 & 0 \\
    -1 & 2 & -1 & 0 \\
    0 & -1 & 2 & -1 \\
    0 & 0 & -1 & 2
  \end{bmatrix},
\end{equation}
with classical reference solution $\mathbf{u}(t) = e^{-At}\,\mathbf{u}_0$.
Since $A$ is real symmetric, the Cartesian decomposition gives $L = A$, $H = 0$,
and the LCHS circuit reduces to a single joint evolution
\begin{equation}\label{eq:joint}
  U(t) = \exp\!\bigl(-it\,\hat{x}\otimes A\bigr),
\end{equation}
with no Trotter splitting required.

The Pauli decomposition of $A$ (at $\alpha = h = 1$) is
\begin{equation}
  A = 2\,I_0I_1 - I_0X_1 - \tfrac{1}{2}(X_0X_1 + Y_0Y_1).
\end{equation}

%% ====================================================================
\section{The LCHS Kernel}
%% ====================================================================

The LCHS integral identity~\cite{das2025} states
\begin{equation}\label{eq:lchs}
  e^{-tA} = \int_{\mathbb{R}} g(k)\,e^{-itkA}\,dk,
\end{equation}
where the near-optimal kernel function is
\begin{equation}\label{eq:kernel}
  g(k) = \frac{f(k)}{1-ik}, \qquad
  f(k) = \frac{e^{2^\beta}}{2\pi\,e^{(1+ik)^\beta}},
  \qquad \beta \in (0,1).
\end{equation}
Equivalently,
\begin{equation}\label{eq:g-combined}
  g(k) = \frac{1}{C_\beta\,e^{(1+ik)^\beta}\,(1-ik)},
  \qquad C_\beta = 2\pi\,e^{-2^\beta}.
\end{equation}

\begin{remark}[Critical constant: $2^\beta$ vs.\ $2\beta$]\label{rem:kernel-bug}
The normalization constant $C_\beta = 2\pi e^{-2^\beta}$ uses
\emph{exponentiation} $2^\beta$, not \emph{multiplication} $2\beta$.
This distinction matters significantly:\ at $\beta = 0.8$,
$2^\beta = 1.741$ while $2\beta = 1.6$, giving
$C_{0.8}^{\text{correct}} = 2\pi e^{-1.741} \approx 1.095$ versus
$C_{0.8}^{\text{wrong}} = 2\pi e^{-1.6} \approx 1.268$.
The error compounds through the coefficient integrals and can invalidate
parameter sweeps performed with the wrong constant.
\end{remark}

%% ====================================================================
\section{Squeezed-Fock State Preparation}
%% ====================================================================

\subsection{Construction}

Following~\cite{das2025}, the post-selection state is a squeezed vacuum
with squeezing parameter $r$:
\begin{equation}\label{eq:phi}
  \phi_r(x) = S(r,0)\phi_0(x) =
  \left(\frac{2}{\sigma^2\pi}\right)^{1/4} e^{-x^2/\sigma^2},
  \qquad \sigma = e^r.
\end{equation}
The requirement $\phi^*(x)\psi(x) = g(x)$ then fixes the preparation state's
position-space wavefunction:
\begin{equation}\label{eq:psi}
  \psi_r(x) = \mathcal{N}\,g(x)\,e^{x^2/\sigma^2},
\end{equation}
which is expanded in a squeezed Fock basis
$|\psi_r\rangle = S(r')\sum_{n=0}^{N-1} C_n|n\rangle$ with coefficients
\begin{equation}\label{eq:Cn}
  C_n = \sqrt{\frac{\sigma}{\sigma'}}\,\frac{1}{\sqrt{2^n n!}}
  \int_{-\infty}^{\infty} H_n\!\left(\frac{x}{\sigma'}\right)
  g(x)\,e^{-\gamma x^2}\,dx,
\end{equation}
where
\begin{equation}\label{eq:gamma}
  \gamma = \frac{1}{\sigma'^2} - \frac{1}{\sigma^2}
  = e^{-2r'} - e^{-2r}, \qquad \sigma' = e^{r'}, \qquad r' < r.
\end{equation}

\subsection{The Fundamental Tension}\label{sec:tension}

Equation~\eqref{eq:psi} reveals a fundamental mathematical tension in the
squeezed-Fock approach.  We require $\psi_r \in L^2(\mathbb{R})$, which demands
\begin{equation}\label{eq:l2-condition}
  \int_{\mathbb{R}} |g(x)|^2\,e^{2x^2/\sigma^2}\,dx < \infty.
\end{equation}

\paragraph{Small $r$ regime ($r \to 0$).}
When $r \to 0$, $\sigma = e^r \to 1$, and the exponential factor
$e^{x^2/\sigma^2} \approx e^{x^2}$ dominates.  Since $|g(x)|$ decays only
super-algebraically (not Gaussian-fast), condition~\eqref{eq:l2-condition}
fails: the integral diverges.  At $r = 0$ exactly, the post-selection state
is the vacuum $\phi_0(x)$, and~\cite[Eq.~10]{das2025} confirms
$\psi_0(x) = \mathcal{N}g(x)e^{x^2} \notin L^2(\mathbb{R})$.

\paragraph{Large $r$ regime ($r \gg 1$).}
When $r$ is large, $\sigma = e^r \gg 1$ and $e^{x^2/\sigma^2} \approx 1$ over
the kernel support $|x| \lesssim 15$, so $\psi_r(x) \approx \mathcal{N}g(x)$
is square-integrable.  However, the post-selection state $|\phi_r\rangle = S(r)|0\rangle$
is a squeezed vacuum whose mean photon number is
\begin{equation}\label{eq:nbar}
  \langle\hat{n}\rangle = \sinh^2(r).
\end{equation}
This grows exponentially with $r$:

\begin{center}
\begin{tabular}{rrcl}
  \toprule
  $r$ & $\sigma = e^r$ & $\langle\hat{n}\rangle = \sinh^2(r)$ & Feasibility at $N = 64$ \\
  \midrule
  $0.003$ & $1.003$ & $9\times10^{-6}$ & Fits trivially; $\sigma \approx 1 \ll$ kernel support \\
  $0.5$ & $1.65$ & $0.27$ & Fits; $\sigma$ still well below kernel support \\
  $1.0$ & $2.72$ & $1.38$ & Fits; $\sigma < 10$ \\
  $2.0$ & $7.39$ & $13.1$ & Fits; $\sigma$ marginally covers support \\
  $3.0$ & $20.1$ & $100$ & Exceeds $N = 64$; state under-resolved \\
  $5.0$ & $148$ & $5{,}500$ & Completely impractical \\
  $8.0$ & $2{,}981$ & $4.5\times10^6$ & Absurd \\
  \bottomrule
\end{tabular}
\end{center}

The condition for $\psi_r$ to be well-approximated is
$\sigma \gg$ kernel support $\approx 15$, requiring $r \gtrsim \ln(15) \approx 2.7$.
But at $r = 2.7$, $\langle\hat{n}\rangle = \sinh^2(2.7) \approx 55$, and a faithful
Fock-space representation needs $N \gg \langle\hat{n}\rangle$.

\begin{proposition}[The squeeze--truncation tradeoff]
For the near-optimal LCHS kernel~\eqref{eq:g-combined} applied to a PDE with
$\|A\| = O(1)$, the squeezed-Fock method requires simultaneously:
\begin{enumerate}[itemsep=2pt]
  \item $\sigma = e^r \gg 15$ for the Gaussian envelope to be flat over the
        kernel support (otherwise $\psi_r \notin L^2$),
  \item $N \gg \sinh^2(r)$ for the Fock truncation to resolve the squeezed
        vacuum (otherwise $|\phi_r\rangle$ is badly approximated),
  \item $r' < r$ with $\gamma > 0$ for the coefficient integral to converge.
\end{enumerate}
These constraints are mutually antagonistic: (1) pushes $r$ large while (2)
demands $N$ grow exponentially with $r$.  The Fock cutoff required for
$\epsilon$-accurate state preparation scales as
$N = \Omega(\sinh^2(\ln(15/\epsilon))) = \Omega(1/\epsilon^2)$,
which is polynomially worse than the $N = O(\log^2(1/\epsilon))$ scaling
for the kernel projection alone (Theorem~2 in~\cite{das2025}).
\end{proposition}

%% ====================================================================
\section{Position-Operator Convention}\label{sec:convention}
%% ====================================================================

The joint Hamiltonian~\eqref{eq:joint} uses $\hat{x}$, the position operator
on the oscillator.  Two conventions appear in the literature:
\begin{equation}\label{eq:conventions}
  \hat{x}_{\text{half}} = \frac{\hat{a}+\hat{a}^\dagger}{2},
  \qquad
  \hat{x}_{\sqrt{2}} = \frac{\hat{a}+\hat{a}^\dagger}{\sqrt{2}}.
\end{equation}
The standard physics convention is $\hat{x}_{\sqrt{2}}$, corresponding to the
canonical commutation relation $[\hat{x},\hat{p}] = i$ with
$\hat{p} = i(\hat{a}^\dagger - \hat{a})/\sqrt{2}$.

The choice of convention directly affects the effective coupling strength in
$\exp(-it\hat{x}\otimes A)$: switching from $\hat{x}_{\sqrt{2}}$ to
$\hat{x}_{\text{half}}$ rescales $\hat{x}$ by a factor of
$\sqrt{2}/2 \approx 0.707$, which is equivalent to running the simulation
at time $t' = 0.707\,t$.  This is not a free parameter---the LCHS integral
identity~\eqref{eq:lchs} is derived for a specific normalization of $\hat{x}$
that must be consistent between the state preparation (where $g(x)$ encodes
position-space amplitudes) and the joint evolution.

\paragraph{Numerical impact.}
At the default parameters ($r = 0.003$, $r' = 0.001$, $\beta = 0.95$,
$N = 64$):
\begin{center}
\begin{tabular}{lcc}
  \toprule
  Convention & Map error (best scale) & Shape fidelity \\
  \midrule
  $\hat{x}_{\text{half}} = (a+a^\dagger)/2$ & $0.337$ & $0.937$ \\
  $\hat{x}_{\sqrt{2}} = (a+a^\dagger)/\sqrt{2}$ & $\mathbf{0.154}$ & $\mathbf{0.988}$ \\
  \bottomrule
\end{tabular}
\end{center}

The $\sqrt{2}$ convention gives significantly better results, consistent with
its being the standard normalization assumed by the LCHS derivation.

\begin{remark}[On empirical scale factors]
Some implementations introduce an empirical \texttt{position\_scale} multiplier
on $\hat{x}$ to improve numerical fit.  This is mathematically unjustified:
any rescaling of $\hat{x}$ changes the effective time or operator norm in
the joint Hamiltonian and breaks the correspondence between the LCHS integral
and the quantum simulation.  The correct approach is to fix the convention
to match the theory derivation.
\end{remark}

%% ====================================================================
\section{Numerical Results}
%% ====================================================================

\subsection{Default Configuration}

\begin{center}
\begin{tabular}{ll}
  \toprule
  Parameter & Default \\
  \midrule
  $\alpha$ & $1.0$ \\
  $h$ & $1.0$ \\
  $t$ & $1.0$ \\
  \texttt{n-fock} & $64$ \\
  \texttt{n-coeff} & $64$ \\
  $r$ (target) & $0.003$ \\
  $r'$ (prime) & $0.001$ \\
  $\beta$ & $0.95$ \\
  \texttt{n-quad} & $300$ \\
  $\hat{x}$ convention & $\sqrt{2}$ \\
  Initial DV state & $|01\rangle$ \\
  \bottomrule
\end{tabular}
\end{center}

\subsection{Baseline Output}

\begin{itemize}[itemsep=2pt]
  \item Operator map relative error (best global scale): $0.1539$
  \item Vector relative error for $|01\rangle$ (best global scale): $0.1084$
  \item Vector shape fidelity: $0.9884$
  \item Eta-calibrated vector error: $0.3855$
  \item Postselection probability: $0.2607$
  \item Best map scale: $+1.2446$
\end{itemize}

\subsection{Metrics Definition}

Given the LCHS map $K$ and reference map $R = e^{-At}$, the best-scale fit is
\begin{equation}
  s_* = \frac{\langle R, K\rangle_F}{\langle R, R\rangle_F},
  \qquad
  \varepsilon_{\text{map}} = \frac{\|K - s_* R\|_F}{\|s_* R\|_F}.
\end{equation}
The eta-calibrated error uses the overlap $\eta = \langle\phi|\psi\rangle$
as the calibration constant:
\begin{equation}
  \varepsilon_{\eta} = \frac{\|\mathbf{u}_{\text{cv}}/\eta - \mathbf{u}_{\text{ref}}\|}
                            {\|\mathbf{u}_{\text{ref}}\|}.
\end{equation}
The best-scale metric allows an optimal complex rescaling (not available on
hardware), while the eta-calibrated metric uses only information available
from the state overlap (physically measurable).

\subsection{Sensitivity to Squeezing Parameter $r$}

Fixing $r' = 0.001$, $\beta = 0.95$, $N = 64$, $\hat{x}_{\sqrt{2}}$:

\begin{center}
\begin{tabular}{rcccc}
  \toprule
  $r$ & $\langle\hat{n}\rangle$ & Map error & Fidelity & Eta error \\
  \midrule
  $0.003$ & $9\times10^{-6}$ & $\mathbf{0.154}$ & $\mathbf{0.988}$ & $0.385$ \\
  $0.1$ & $0.01$ & $0.245$ & $0.969$ & $0.501$ \\
  $0.3$ & $0.09$ & $0.373$ & $0.925$ & $0.666$ \\
  $0.5$ & $0.27$ & $0.471$ & $0.875$ & $0.805$ \\
  $0.693$ & $0.57$ & $0.556$ & $0.822$ & $0.937$ \\
  $1.0$ & $1.38$ & $0.681$ & $0.734$ & $1.154$ \\
  $1.5$ & $4.53$ & $0.843$ & $0.616$ & $1.465$ \\
  $2.0$ & $13.1$ & $0.935$ & $0.553$ & $1.654$ \\
  \bottomrule
\end{tabular}
\end{center}

The error increases \emph{monotonically} with $r$.  This is the opposite of
the theoretical expectation (larger $r$ should give better $L^2$ approximation
to $\psi_r$) and reflects the dominance of the Fock truncation error:
the squeezed vacuum $S(r)|0\rangle$ becomes increasingly poorly represented
in the finite Fock space as $r$ grows, and this truncation error overwhelms
the improvement in the kernel overlap.

\subsection{Sensitivity to Kernel Parameter $\beta$}

Fixing $r = 0.003$, $r' = 0.001$, $N = 64$, $\hat{x}_{\sqrt{2}}$:

\begin{center}
\begin{tabular}{cccc}
  \toprule
  $\beta$ & Map error & Fidelity & Eta error \\
  \midrule
  $0.3$ & $0.354$ & $0.935$ & $0.805$ \\
  $0.5$ & $0.294$ & $0.956$ & $0.669$ \\
  $0.7$ & $0.233$ & $0.973$ & $0.541$ \\
  $0.8$ & $0.202$ & $0.980$ & $0.478$ \\
  $0.9$ & $0.170$ & $0.986$ & $0.416$ \\
  $0.95$ & $\mathbf{0.154}$ & $\mathbf{0.988}$ & $\mathbf{0.385}$ \\
  $0.99$ & $0.141$ & $0.990$ & $0.361$ \\
  \bottomrule
\end{tabular}
\end{center}

Higher $\beta$ consistently improves accuracy.  This is expected: larger
$\beta$ narrows the kernel $g(k)$, concentrating its support closer to
$k = 0$ where the finite Fock space can better resolve the overlap integral.
The improvement is smooth and monotonic, approaching a floor near
$\beta \to 1$.

\subsection{Fock Dimension Convergence}

Fixing $r = 0.003$, $r' = 0.001$, $\beta = 0.95$, $\hat{x}_{\sqrt{2}}$:

\begin{center}
\begin{tabular}{rcc}
  \toprule
  $N$ & Map error & Fidelity \\
  \midrule
  $16$ & $0.1539$ & $0.9884$ \\
  $32$ & $0.1539$ & $0.9884$ \\
  $64$ & $0.1539$ & $0.9884$ \\
  $128$ & $0.1539$ & $0.9884$ \\
  \bottomrule
\end{tabular}
\end{center}

The error is completely saturated: increasing $N$ beyond 16 provides no
improvement.  At $r \approx 0$, the squeezed states are essentially
unsqueezed ($S(0.003) \approx I$), so the Fock-space representation is
trivially exact.  The residual 15.4\% error comes entirely from the
kernel approximation---the finite squeezed-Fock expansion simply cannot
represent $g(x)$ well in the near-vacuum limit.

\subsection{Trajectory}

At the default parameters, the best-scale error along the trajectory is:

\begin{center}
\begin{tabular}{cccc}
  \toprule
  Time $t$ & Best-scale error & Fidelity & Post.\ probability \\
  \midrule
  $0.000$ & $0.000$ & $1.000$ & $0.786$ \\
  $0.500$ & $0.183$ & $0.968$ & $0.390$ \\
  $0.740$ & $0.143$ & $0.980$ & $0.308$ \\
  $1.000$ & $0.108$ & $0.988$ & $0.261$ \\
  \bottomrule
\end{tabular}
\end{center}

The error peaks near $t \approx 0.36$ (maximum best-scale error $0.198$)
and decreases toward $t = 1$, consistent with the quadrature accuracy
being most challenged at intermediate times.

%% ====================================================================
\section{The Kernel Parameter $\beta$: Why $\beta \in (0,1)$ and Cannot Equal~1}\label{sec:beta}
%% ====================================================================

The sensitivity data in Section~\ref{sec:convention} shows that accuracy
improves monotonically as $\beta \to 1$.  A natural question is: why not
set $\beta = 1$?  We show that $\beta = 1$ is a singular limit at which
the kernel ceases to be integrable, and that the approach to this limit
involves a precise tradeoff between kernel concentration and support
growth.

\subsection{Asymptotic decay of $g(k)$ for large $|k|$}

For $|k| \gg 1$, write $1 + ik = ik(1 - i/k)$.  Using
$ik = |k|\,e^{i\pi/2}$ (for $k > 0$), the principal branch gives
\begin{equation}\label{eq:asymp-base}
  (1+ik)^\beta \;=\; |k|^\beta\,e^{i\beta\pi/2}\,\bigl(1 + O(1/k)\bigr),
  \qquad k \to +\infty.
\end{equation}
Taking real and imaginary parts:
\begin{align}
  \operatorname{Re}\!\bigl[(1+ik)^\beta\bigr]
  &\;=\; k^\beta\cos\!\bigl(\tfrac{\beta\pi}{2}\bigr) + O(k^{\beta-1}),
  \label{eq:re-asymp} \\
  \operatorname{Im}\!\bigl[(1+ik)^\beta\bigr]
  &\;=\; k^\beta\sin\!\bigl(\tfrac{\beta\pi}{2}\bigr) + O(k^{\beta-1}).
  \label{eq:im-asymp}
\end{align}
From Eq.~\eqref{eq:g-combined}, $|g(k)| = |C_\beta|^{-1}\,
\exp\!\bigl(-\operatorname{Re}[(1+ik)^\beta]\bigr)\,/\,\sqrt{1+k^2}$, so
\begin{equation}\label{eq:g-decay}
  \boxed{%
  |g(k)| \;\sim\; \frac{1}{|C_\beta|\,|k|}\,
  \exp\!\Bigl(-|k|^\beta\cos\!\bigl(\tfrac{\beta\pi}{2}\bigr)\Bigr),
  \qquad |k| \to \infty.}
\end{equation}

\subsection{The $\beta < 1$ regime: super-algebraic decay}

For any $\beta \in (0,1)$, $\cos(\beta\pi/2) > 0$.  The exponential
factor $\exp(-|k|^\beta\cos(\beta\pi/2))$ provides
\emph{super-algebraic} (stretched-exponential) decay, which dominates
the algebraic $1/|k|$ prefactor.  In particular:
\begin{equation}\label{eq:g-integrable}
  \int_{\mathbb{R}} |g(k)|\,dk \;<\; \infty
  \qquad \text{for all } \beta \in (0,1).
\end{equation}
This is the fundamental requirement for the LCHS identity~\eqref{eq:lchs}
to hold as an absolutely convergent integral.

\subsection{The $\beta = 1$ singularity}

At $\beta = 1$:
\begin{equation}\label{eq:cos-zero}
  \cos\!\bigl(\tfrac{\pi}{2}\bigr) = 0.
\end{equation}
The exponential decay factor vanishes, and Eq.~\eqref{eq:g-decay} reduces to
\begin{equation}\label{eq:g-beta1}
  |g(k)| \;\sim\; \frac{1}{|C_1|\,|k|}, \qquad |k| \to \infty,
  \qquad (\beta = 1).
\end{equation}
Since $\int_1^\infty dk/k = \infty$, the kernel is \emph{not} absolutely
integrable at $\beta = 1$:
\begin{equation}\label{eq:diverge}
  \int_{\mathbb{R}} |g(k)|\,dk \;=\; \infty
  \qquad \text{at } \beta = 1.
\end{equation}
The LCHS integral~\eqref{eq:lchs} therefore does not converge absolutely,
and the kernel cannot be used as a probability amplitude for quantum
state preparation (which requires $g \in L^1(\mathbb{R})$).

\begin{remark}[Exact form at $\beta = 1$]
At $\beta = 1$, $(1+ik)^\beta = 1 + ik$ exactly, so
$f(k) = e^{2}/(2\pi\,e^{1+ik}) = e/(2\pi)\,e^{-ik}$ and
$g(k) = e/(2\pi(1-ik))\,e^{-ik}$.
This has $|g(k)| = e/(2\pi\sqrt{1+k^2})$, confirming the $1/|k|$ decay.
The Fourier transform of $1/(1-ik)$ is a one-sided exponential, so $g(k)$
at $\beta = 1$ is essentially a phase-shifted Lorentzian---integrable
but only barely, and not absolutely integrable in the $L^1$ sense needed
for LCHS.
\end{remark}

\subsection{Effective kernel support and $L^1$ norm}

Define the effective support half-width $k_{\mathrm{eff}}$ as the scale
at which the exponential decay factor equals $1/e$:
\begin{equation}\label{eq:keff}
  k_{\mathrm{eff}}^\beta \cos\!\bigl(\tfrac{\beta\pi}{2}\bigr) = 1
  \qquad\Longrightarrow\qquad
  k_{\mathrm{eff}} = \left(\frac{1}{\cos(\beta\pi/2)}\right)^{1/\beta}.
\end{equation}
As $\beta \to 1^-$, $\cos(\beta\pi/2) \to 0^+$, and $k_{\mathrm{eff}}$
diverges:
\begin{equation}\label{eq:keff-diverge}
  k_{\mathrm{eff}} \to \infty \quad \text{as } \beta \to 1^-.
\end{equation}

Numerical evaluation of $k_{\mathrm{eff}}$ and the $L^1$ norm
$\|g\|_1 = \int |g(k)|\,dk$:

\begin{center}
\begin{tabular}{rccl}
  \toprule
  $\beta$ & $\cos(\beta\pi/2)$ & $k_{\mathrm{eff}}$ & $\|g\|_1$ \\
  \midrule
  $0.30$ & $0.891$ & $1.47$ & $1.03$ \\
  $0.50$ & $0.707$ & $2.00$ & $1.10$ \\
  $0.70$ & $0.454$ & $3.52$ & $1.28$ \\
  $0.80$ & $0.309$ & $5.82$ & $1.48$ \\
  $0.90$ & $0.156$ & $13.4$ & $1.87$ \\
  $0.95$ & $0.0785$ & $32.1$ & $2.33$ \\
  $0.99$ & $0.0157$ & $66.4$ & $3.73$ \\
  $0.999$ & $0.00157$ & --- & $4.94$ \\
  $1.00$ & $0$ & $\infty$ & $\infty$ \\
  \bottomrule
\end{tabular}
\end{center}

\subsection{The $\beta$--squeezing tradeoff for LCHS}

The LCHS quantum circuit encodes $g(k)$ in the position-space
wavefunction of a bosonic mode.  A Fock space of dimension $N$ can
faithfully represent wavefunctions whose position-space support fits
within $|x| \lesssim \sqrt{2N}$ (the range of the highest Fock-state
eigenfunction).  This imposes the constraint
\begin{equation}\label{eq:fock-support}
  k_{\mathrm{eff}} \;\lesssim\; \sqrt{2N}.
\end{equation}
Combining with Eq.~\eqref{eq:keff}:
\begin{equation}\label{eq:beta-N-constraint}
  \left(\frac{1}{\cos(\beta\pi/2)}\right)^{1/\beta} \lesssim \sqrt{2N}
  \qquad\Longleftrightarrow\qquad
  \cos\!\bigl(\tfrac{\beta\pi}{2}\bigr)
  \;\gtrsim\; (2N)^{-\beta/2}.
\end{equation}
At $\beta = 0.95$ with $N = 64$: $k_{\mathrm{eff}} \approx 32$ and
$\sqrt{2 \times 64} \approx 11.3$, so we already violate
condition~\eqref{eq:fock-support}.  This explains why higher $\beta$
improves accuracy (kernel concentrates near $k = 0$, where the finite
Fock representation is accurate) but with diminishing returns (the
tails beyond $\sqrt{2N}$ contribute growing $L^1$ mass that cannot
be captured).

At $\beta = 1$, $k_{\mathrm{eff}} = \infty$ and no finite $N$ suffices:
the kernel has infinite support, infinite $L^1$ norm, and cannot be
represented in any finite-dimensional Fock space.  This is the
mathematical reason $\beta = 1$ is excluded.

\subsection{Physical interpretation}

In the LCHS framework, $\beta$ controls the shape of the spectral
filter that converts unitary evolution $e^{-itkA}$ into dissipative
evolution $e^{-tA}$.  The kernel $g(k)$ is the filter function whose
Fourier structure must satisfy:
\begin{enumerate}[itemsep=2pt]
  \item \textbf{Normalization}: $\int g(k)\,dk = 1$ (the identity
        $e^{-tA}$ requires unit integral).
  \item \textbf{Absolute convergence}: $\int |g(k)|\,dk < \infty$
        (the integral must converge as an operator norm bound).
  \item \textbf{Finite encoding}: $g(k)$ must be representable in
        a finite-dimensional quantum register.
\end{enumerate}

Condition~2 fails at $\beta = 1$ (Eq.~\ref{eq:diverge}).
Even if one relaxes to conditional convergence, condition~3 fails
because the support $k_{\mathrm{eff}} = \infty$ exceeds any finite
Fock space.

The parameter $\beta$ thus lives in the open interval $(0,1)$, with
$\beta \to 1^-$ giving the best spectral concentration but at the
cost of growing $k_{\mathrm{eff}}$ and $\|g\|_1$.  The optimal $\beta$
for a given Fock dimension $N$ is the largest value satisfying
Eq.~\eqref{eq:fock-support}, roughly
\begin{equation}\label{eq:beta-opt}
  \beta_{\max}(N) \;=\; \max\Bigl\{\beta \in (0,1)
  \;:\; \bigl(\cos(\beta\pi/2)\bigr)^{-1/\beta} \le \sqrt{2N}\Bigr\}.
\end{equation}

%% ====================================================================
\section{Why Near-Zero Squeezing Works Best}
%% ====================================================================

The seemingly paradoxical result that $r \approx 0$ outperforms $r > 0$
can be understood as follows.

\paragraph{1.\ The coefficient integral is well-behaved at small $r$.}
From Eq.~\eqref{eq:Cn}, the Gaussian envelope $e^{-\gamma x^2}$ with
$\gamma = e^{-2r'} - e^{-2r}$ acts as a convergence factor for the
integral.  When both $r$ and $r'$ are small, $\gamma$ is small but positive
(e.g., $\gamma \approx 0.004$ at the defaults), so the envelope is broad
and the integral samples $g(x)$ over a wide range.  The Hermite polynomials
$H_n(x/\sigma')$ with $\sigma' \approx 1$ provide an efficient basis for
capturing the kernel's structure near $x = 0$.

\paragraph{2.\ The squeezed operators are near-identity.}
At $r \approx 0$, $S(r) \approx I$, so the state preparation simplifies to
$|\psi\rangle \approx \sum_n C_n |n\rangle$ and
$|\phi\rangle \approx |0\rangle$.
The Fock space trivially contains both states, and there is no truncation
error from squeezing.

\paragraph{3.\ The residual error is from the kernel projection.}
The 15.4\% best-scale error at $r \approx 0$ represents the intrinsic
limitation of the squeezed-Fock expansion in the near-vacuum limit:
$\psi_0(x) = \mathcal{N}g(x)e^{x^2}$ is not normalizable, so any finite
$r > 0$ gives a $\psi_r$ that deviates from the ideal kernel.  The
expansion coefficients $\{C_n\}$ can only approximate $g(x)$ to the extent
that the Gaussian-weighted Hermite projection captures the kernel's
non-Gaussian features.

\paragraph{4.\ Increasing $r$ introduces truncation error faster than it
improves the kernel approximation.}
The improvement in the $L^2$ norm of $\psi_r$ from increasing $r$ is
gradual (the Gaussian envelope $e^{x^2/\sigma^2}$ flattens slowly), while
the Fock truncation error from representing $S(r)|0\rangle$ in finite
dimension grows as $\sinh^2(r)$.  The net effect is always negative at
practical Fock dimensions.

%% ====================================================================
\section{Comparison with Position-Eigenbasis Method}
%% ====================================================================

The position-eigenbasis method (implemented in \texttt{claude\_heat1d\_qutip.py})
avoids the squeezed-state construction entirely by encoding $g(x)$ directly
in the eigenbasis of $\hat{x}$.  The comparison is stark:

\begin{center}
\begin{tabular}{lccc}
  \toprule
  Method & Best params & Map error & Fidelity \\
  \midrule
  Squeezed-Fock ($r \approx 0$, $\beta = 0.95$, $N=64$)
    & This note & $15.4\%$ & $98.8\%$ \\
  Position-eigenbasis ($\beta = 0.8$, $N = 100$)
    & \cite{claude_xbasis} & $4.8\%$ & --- \\
  Position-eigenbasis ($\beta = 0.8$, $N = 300$)
    & \cite{claude_xbasis} & $0.2\%$ & --- \\
  \bottomrule
\end{tabular}
\end{center}

The position-eigenbasis method achieves two orders of magnitude better
accuracy because:
\begin{enumerate}[itemsep=3pt]
  \item It avoids the problematic $e^{x^2/\sigma^2}$ factor entirely.
  \item The $\hat{x}$-eigenvalues provide a natural quadrature grid whose
        range grows as $\sqrt{2N}$, automatically covering the kernel support.
  \item The error decreases monotonically with $N$ (pure quadrature convergence).
  \item No squeezing parameters need tuning.
\end{enumerate}

The tradeoff is that the position-eigenbasis encoding does not correspond
to a simple physical state-preparation circuit---it serves as an idealized
numerical reference.

%% ====================================================================
\section{Kernel Constant Bug in Prior Work}\label{sec:kernel-bug}
%% ====================================================================

An earlier implementation (\texttt{OLD\_FILES\_2ND/heat1d\_lchs\_cv\_dv.py})
used the kernel normalization
\begin{equation}
  C_\beta^{\text{wrong}} = 2\pi\,e^{-2\beta}
  \qquad\text{(multiplication)},
\end{equation}
instead of the correct
\begin{equation}
  C_\beta^{\text{correct}} = 2\pi\,e^{-2^\beta}
  \qquad\text{(exponentiation)}.
\end{equation}
A parameter sensitivity sweep was performed with this incorrect kernel,
finding optimal parameters $r \approx 0.48$, $r' \approx 0.40$,
$\beta \approx 0.35$ at $N = 32$.  These parameters achieve
$\sim\!19\%$ PDE error with the wrong kernel but are not valid for the
correct one.

The distinction matters because $2^\beta \neq 2\beta$ for all
$\beta \in (0,1)$ except $\beta \approx 0.6346$ (the fixed point).
At $\beta = 0.35$: $2^{0.35} = 1.275$ vs.\ $2\times0.35 = 0.70$,
a 82\% relative difference in the exponent.

%% ====================================================================
\section{Bosonic-Qiskit Circuit Implementation}\label{sec:bosonic}
%% ====================================================================

The QuTiP reference (\texttt{clacod\_heat1d\_qutip.py}) computes the joint
evolution $\exp(-it\hat{x}\otimes A)$ exactly via matrix exponentiation.
To validate the actual quantum circuit, we translate the same LCHS protocol
into a Trotterized circuit on \texttt{bosonic-qiskit}, implemented in
\texttt{clacod\_heat1d\_bosonic.py}.

\subsection{Displacement Convention Derivation}

The LCHS derivation assumes the standard physics position operator
$\hat{x} = (a+a^\dagger)/\sqrt{2}$, satisfying $[\hat{x},\hat{p}]=i$.
The \texttt{bosonic-qiskit} displacement gate implements the standard
Glauber operator
\begin{equation}
  D(\alpha) = \exp(\alpha\, a^\dagger - \alpha^*\, a).
\end{equation}
To find the displacement amplitude $\alpha$ that generates
$\exp(-i\lambda\hat{x})$ for real $\lambda$:
\begin{align}
  \exp(-i\lambda\hat{x})
  &= \exp\!\left(\frac{-i\lambda}{\sqrt{2}}(a + a^\dagger)\right) \notag\\
  &= \exp\!\left(\frac{-i\lambda}{\sqrt{2}}\,a^\dagger
    + \frac{-i\lambda}{\sqrt{2}}\,a\right).
\end{align}
Matching $\alpha\,a^\dagger - \alpha^*\,a
= \tfrac{-i\lambda}{\sqrt{2}}\,a^\dagger + \tfrac{-i\lambda}{\sqrt{2}}\,a$
gives
\begin{equation}\label{eq:disp-alpha}
  \boxed{\alpha = \frac{-i\lambda}{\sqrt{2}},}
\end{equation}
which is consistent since $(-i\lambda/\sqrt{2})^* = i\lambda/\sqrt{2}$
and $-\alpha^* = -i\lambda/\sqrt{2}$.

\subsection{Conditional Displacement for Pauli Terms}

The \texttt{bosonic-qiskit} conditional displacement
$\texttt{cv\_c\_d}(\alpha, \text{qubit})$ applies
\begin{equation}
  \mathrm{CD}(\alpha) = |0\rangle\!\langle 0|\otimes D(\alpha)
  + |1\rangle\!\langle 1|\otimes D(-\alpha).
\end{equation}
For the operator $\exp(-i\lambda\,\hat{x}\otimes Z)$, the qubit $|0\rangle$
branch needs $\exp(-i\lambda\hat{x}) = D(-i\lambda/\sqrt{2})$ and
the $|1\rangle$ branch needs $\exp(+i\lambda\hat{x}) = D(+i\lambda/\sqrt{2})$.
This matches $\mathrm{CD}(\alpha)$ with $\alpha = -i\lambda/\sqrt{2}$.

For a general Pauli string $P$, we diagonalize it into a single-qubit $Z$
via Clifford basis changes and a CNOT parity-compute circuit:
\begin{itemize}[itemsep=2pt]
  \item $X = H Z H$, so $\exp(-i\lambda\hat{x}\otimes X)$ uses Hadamard
        conjugation on the target qubit.
  \item $X\otimes X = (H\otimes H)(Z\otimes Z)(H\otimes H)$ and
        $Z\otimes Z = \mathrm{CNOT}\cdot(I\otimes Z)\cdot\mathrm{CNOT}$.
  \item $Y\otimes Y$ uses the additional basis change
        $Y = R_z(\pi/2)\,H\,Z\,H\,R_z(-\pi/2)$ to map $Y$ eigenstates
        to $Z$ eigenstates.
\end{itemize}

\subsection{Trotterized Circuit per Step}

For the heat equation $A = \tfrac{\alpha}{h^2}(2\cdot II - IX
- \tfrac{1}{2}XX - \tfrac{1}{2}YY)$, each first-order Trotter step
at time slice $\delta t = t/n$ implements:
\begin{equation}
  \mathcal{U}(\delta t)
  \approx
  \prod_{k \in \{II,\, IX,\, XX,\, YY\}}
  \exp\!\left(-i\delta t\, c_k\, \hat{x} \otimes P_k\right),
\end{equation}
with displacement amplitude $\alpha_k = -i\,\delta t\, c_k / \sqrt{2}$
for each Pauli term.  The per-step gate decomposition is:

\begin{center}
\begin{tabular}{llccccc}
  \toprule
  Pauli term & Circuit pattern & D & CD & H & CNOT & $R_z$ \\
  \midrule
  $II$ & $D(\alpha_{II})$ & 1 & --- & --- & --- & --- \\
  $IX$ & $H$-$\mathrm{CD}$-$H$ & --- & 1 & 2 & --- & --- \\
  $XX$ & $HH$-$\mathrm{CX}$-$\mathrm{CD}$-$\mathrm{CX}$-$HH$ & --- & 1 & 4 & 2 & --- \\
  $YY$ & $R_z R_z$-$HH$-$\mathrm{CX}$-$\mathrm{CD}$-$\mathrm{CX}$-$HH$-$R_z R_z$
       & --- & 1 & 4 & 2 & 4 \\
  \midrule
  \textbf{Total per step} & & \textbf{1} & \textbf{3} & \textbf{10}
  & \textbf{4} & \textbf{4} \\
  \bottomrule
\end{tabular}
\end{center}

\subsection{CV State Injection}

Gate-level preparation of the non-Gaussian state
$|\psi\rangle = S(r')\sum_n C_n|n\rangle$ is a hard problem.
We bypass it using \texttt{cv\_initialize}, which directly sets the
Fock amplitudes $\{C_n\}$ followed by squeezing $S(r')$.
This is a \emph{simulator-only instruction}, not a physical gate;
gate-based state preparation will be addressed in future work.

The post-selection projection $\langle\varphi_r| = \langle 0|S^\dagger(r)$
is implemented physically: apply $S(-r)$ to the oscillator mode and
extract the Fock $|0\rangle$ component of the DV state.

\subsection{Circuit Resource Statistics}

For the 1D heat equation with $N=4$ grid points ($m=2$ qubits, 1 qumode),
the total physical gate counts scale linearly with Trotter steps~$n$:

\begin{center}
\begin{tabular}{lcc}
  \toprule
  Gate type & Per step & Total ($n$ steps) \\
  \midrule
  $D$ (displacement, CV) & 1 & $n$ \\
  CD (conditional displacement, hybrid) & 3 & $3n$ \\
  $S$ (squeeze, CV) & --- & 2 \\
  $H$ (Hadamard, DV) & 10 & $10n$ \\
  CNOT (DV) & 4 & $4n$ \\
  $R_z$ (DV) & 4 & $4n$ \\
  \midrule
  \textbf{Physical CV gates} & 4 & $4n + 2$ \\
  \textbf{Physical DV gates} & 18 & $18n$ \\
  \textbf{Total physical} & 22 & $22n + 2$ \\
  \bottomrule
\end{tabular}
\end{center}

\noindent
The system uses 2 DV qubits + 1 CV qumode (Fock dimension $N_F$,
encoded in $\log_2 N_F$ mode qubits for simulation).
The circuit depth estimate is $22n + 2$.

\begin{center}
\begin{tabular}{rccc}
  \toprule
  $n$ (Trotter steps) & Physical gates & Hybrid CD gates & Depth \\
  \midrule
  100 & 2{,}202 & 300 & 2{,}202 \\
  200 & 4{,}402 & 600 & 4{,}402 \\
  \bottomrule
\end{tabular}
\end{center}

\subsection{Numerical Validation: Bosonic Circuit vs.\ QuTiP Reference}

The bosonic-qiskit circuit is compared against both the QuTiP exact result
(no Trotterization) and the classical reference $\exp(-At)\mathbf{u}_0$.
At default parameters ($r=0.003$, $r'=0.001$, $\beta=0.95$, $N_F=64$,
initial state $|01\rangle$):

\begin{center}
\begin{tabular}{lccc}
  \toprule
  Comparison & Rel.\ error (best scale) & Shape fidelity & Scale factor \\
  \midrule
  Bosonic vs.\ Classical & $10.70\%$ & $0.9887$ & $1.201$ \\
  Bosonic vs.\ QuTiP & $\mathbf{0.16\%}$ & $\mathbf{0.99999}$ & $0.999$ \\
  QuTiP vs.\ Classical & $10.84\%$ & $0.9886$ & $1.200$ \\
  \bottomrule
\end{tabular}
\end{center}

\paragraph{Key observations.}
\begin{enumerate}[itemsep=2pt]
  \item The bosonic circuit matches the QuTiP exact LCHS result to $0.16\%$,
        confirming that Trotterization with $n=200$ steps introduces negligible error.
  \item The $\sim\!10.7\%$ gap to the classical solution is identical in both
        QuTiP and bosonic implementations---this is the intrinsic LCHS kernel
        approximation error, not a circuit artifact.
  \item The output vectors from all three methods agree in direction
        (fidelity $> 0.988$) and differ only by a global scale factor
        $\approx 1.20$, consistent with the postselection normalization.
  \item The postselection probability is $p_{\mathrm{succ}} \approx 0.260$,
        matching the QuTiP prediction.
\end{enumerate}

\subsection{Trotter Convergence}

At fixed CV parameters, the Trotter error (bosonic vs.\ QuTiP) converges
as expected:

\begin{center}
\begin{tabular}{rcc}
  \toprule
  $n$ (steps) & Trotter error (vs.\ QuTiP) & Fidelity (vs.\ QuTiP) \\
  \midrule
  100 & $0.15\%$ & $0.999998$ \\
  200 & $0.16\%$ & $0.999997$ \\
  \bottomrule
\end{tabular}
\end{center}

Both values are well below the $10.7\%$ kernel approximation floor, confirming
that the Trotterized circuit faithfully reproduces the continuous LCHS integral.

%% ====================================================================
\section{Reproducibility}
%% ====================================================================

\subsection{QuTiP reference (exact LCHS)}
\begin{lstlisting}[language=bash]
python3 clacod_heat1d_qutip.py
\end{lstlisting}

\subsection{Bosonic-qiskit circuit}
\begin{lstlisting}[language=bash]
python3 clacod_heat1d_bosonic.py \
  --max-fock-level 64 \
  --n-trotter-steps 200 \
  --r-target 0.003 --r-prime 0.001 \
  --beta 0.95 --n-coeff 64
\end{lstlisting}

\subsection{Trajectory with JSON output}
\begin{lstlisting}[language=bash]
python3 clacod_heat1d_qutip.py \
  --trajectory-steps 50 \
  --output-json results_bosonic/clacod_metrics.json
\end{lstlisting}

\subsection{Parameter sweep}
\begin{lstlisting}[language=bash]
python3 clacod_heat1d_3param_sweep.py \
  --method grid+nm --n-fock 64
\end{lstlisting}

%% ====================================================================
\section{Limitations and Conclusions}
%% ====================================================================

\begin{enumerate}[itemsep=3pt]
  \item \textbf{The squeezed-Fock method is fundamentally limited at practical
        Fock dimensions.}
        The best achievable accuracy is $\sim\!15\%$ best-scale error ($\sim\!39\%$
        eta-calibrated) at $N = 64$ with near-zero squeezing.  Increasing $N$ does
        not help because the error is from the kernel projection, not the Fock
        truncation.

  \item \textbf{Larger squeezing makes results worse, not better.}
        Despite the theoretical motivation for large $r$ (flattening the
        Gaussian envelope), the Fock truncation error from representing
        $S(r)|0\rangle$ grows faster than the kernel approximation improves.

  \item \textbf{The position-operator convention matters.}
        The $\hat{x}_{\sqrt{2}} = (a+a^\dagger)/\sqrt{2}$ convention gives
        $2.2\times$ better map error than $\hat{x}_{\text{half}} = (a+a^\dagger)/2$,
        consistent with the LCHS derivation assuming the standard physics convention.

  \item \textbf{Higher $\beta$ helps.}
        Values of $\beta$ close to 1 concentrate the kernel near $k = 0$,
        improving the finite-dimensional projection quality.

  \item \textbf{The bosonic-qiskit circuit validates the QuTiP reference.}
        With Trotterized conditional displacements, the circuit reproduces
        the exact LCHS map to $0.16\%$ relative error (at $n=200$ Trotter steps),
        confirming that the Pauli decomposition, basis-change circuits, and
        displacement conventions are all correct.

  \item \textbf{CV state injection bypasses the hard state-prep problem.}
        Using \texttt{cv\_initialize} allows validating the Hamiltonian evolution
        circuit independently of gate-level state preparation.
        Gate-based preparation of $S(r')\sum C_n|n\rangle$ remains open for
        future work (e.g., SNAP gates, optimal control pulses, or sequential
        number-selective rotations).

  \item \textbf{The position-eigenbasis method is strictly superior for numerical
        validation.}
        It achieves $0.2\%$ error at $N = 300$ with no squeezing parameters to tune,
        versus $15.4\%$ for the best squeezed-Fock configuration.

  \item \textbf{The squeezed-Fock method's value is in its physical realizability.}
        Despite the numerical limitations, it maps to a concrete circuit:
        prepare $S(r')\sum C_n|n\rangle$, evolve under $\hat{x}\otimes A$
        via conditional displacements, and post-select on $\langle 0|S^\dagger(r)$.
        The circuit resource cost scales as $22n + 2$ physical gates for
        $n$ Trotter steps on the 4-point heat equation (2 qubits, 1 qumode).
\end{enumerate}

\begin{thebibliography}{9}
\bibitem{das2025}
E.~R.~Das, M.~Zheng, R.~Dutta, A.~Li, and Y.~Liu,
``Hybrid Oscillator-Qubit Linear Combination of Hamiltonian Simulation,''
2025.

\bibitem{an2023}
D.~An, J.-P.~Liu, and L.~Lin,
``Linear combination of Hamiltonian simulation for non-unitary dynamics
with optimal state preparation cost,''
\textit{Phys.\ Rev.\ Lett.}, vol.~131, 150603, 2023.

\bibitem{an2023b}
D.~An, A.~M.~Childs, and L.~Lin,
``Quantum algorithm for linear non-unitary dynamics with near-optimal
dependence on all parameters,''
\textit{arXiv:2312.03916}, 2023.

\bibitem{claude_xbasis}
Position-eigenbasis implementation:
\texttt{claude\_heat1d\_qutip.py} in this repository.
\end{thebibliography}

\end{document}
